% ECONOMICS_PROBLEM: {"id": "Q54-223", "title": "Climate tipping points and nonlinear damages", "jel_code": "Q54", "source_id": "OP-0112"}
% !TeX program = xelatex
\documentclass[11pt,a4paper]{article}
\usepackage[margin=23mm]{geometry}
\usepackage{fontspec}
\setmainfont{Latin Modern Roman}
\usepackage{amsmath,amssymb,mathtools}
\usepackage{xcolor,enumitem,booktabs,longtable,array,xurl,hyperref}
\definecolor{muted}{HTML}{53636C}
\hypersetup{colorlinks=true,urlcolor=blue,linkcolor=blue}
\setlength{\parindent}{0pt}
\setlength{\parskip}{5pt}
\newcommand{\R}{\mathbb R}
\newcommand{\E}{\mathbb E}
\newcommand{\N}{\mathbb N}
\newcommand{\ind}{\mathbf 1}
\newcommand{\Var}{\operatorname{Var}}
\newcommand{\Cov}{\operatorname{Cov}}
\newcommand{\supp}{\operatorname{supp}}
\DeclareMathOperator*{\argmin}{arg\,min}
\DeclareMathOperator*{\argmax}{arg\,max}
\newcommand{\entrymeta}[3]{{\sffamily\small\color{muted}\textbf{\texttt{#1}}\quad|\quad #2\par #3\par}}
\newcommand{\Problemref}[1]{\texttt{#1}}
\newcommand{\JELref}[2]{\texttt{#2} (JEL \texttt{#1})}
\begin{document}
\section*{Climate tipping points and nonlinear damages}
\textbf{Problem ID:} \texttt{Q54-223}\quad \textbf{JEL:} \texttt{Q54}\par
% BEGIN ECONOMICS STATEMENT
\label{problem:OP-0112}
\label{atlas:prob:C2}

\noindent\textbf{Original identifier:} C2 (atlas item 112). \textbf{Field:} Environmental and climate economics. \textbf{Classification:} editorial research family with established restricted solutions; current open status not certified.

\subsection*{Definitions and assumptions}

Consider dates $t=0,\ldots,H$, climate-economic state $x_t$, mitigation/adaptation action $a_t\in A(x_t,J_t)$ and tipping indicators $J_t\in\{0,1\}^K$ for $K$ specified physical elements. Irreversibility over this horizon means $J_{t+1,k}\ge J_{t,k}$; reversibility or hysteresis needs another state representation. A transition kernel $Q(dx',dj'\mid x,j,a)$ contains temperature-dependent hazards and interactions. Per-period welfare $U(x,j,a)$ incorporates consumption, ecological services and action costs in a declared utility scale; discount factor $\beta\in(0,1]$ and terminal value $G$ are supplied. A policy $\pi$ is a measurable mapping from the observed history into feasible actions.

\subsection*{Formal research question}

For transition kernels compatible with physical evidence and stated economic propagation restrictions, characterize optimal mitigation and tipping-risk values. A time-consistent robust formulation, when uncertainty is rectangular across histories, has Bellman recursion
\[
V_t(x,j)=\sup_{a\in A(x,j)}
\left\{U(x,j,a)+\beta\inf_{Q\in\mathcal Q_t(x,j,a)}
\mathbb E_Q[V_{t+1}(x',j')]\right\},\quad V_H=G.
\]
Here $\mathcal Q_t$ is an explicitly specified nonempty transition set. Establish existence and approximation-error bounds under appropriate boundedness/continuity conditions; identify which hazard, damage and interaction restrictions change the optimal action robustly. Monetize welfare changes using a positive marginal consumption value in a stated real-currency year, and evaluate how adding migration or capital destruction changes the range of policy conclusions.

\subsection*{Interpretation and scope}

A threshold in a physical state is not automatically a discontinuity in aggregate monetary damages. Economic effects depend on exposure, adaptation and substitution, and a probability of tipping differs from the distribution of transition times. Rectangularity permits nature's conditional choices to vary by history; a single fixed but unknown physical model defines a different ambiguity problem and can invalidate this Bellman recursion. Bayesian learning requires a prior and observation process. A finite-horizon irreversibility approximation must not be presented as a universal physical fact. Sensitivity to discounting and regional welfare weights should be reported separately from estimated hazard uncertainty.

\subsection*{Source audit and status}

[S1] identifies physical tipping elements; [S2] studies optimal policy in a specified tipping IAM. The input's Cai--Judd--Lontzek (2016) does not uniquely match the relevant publication: the verified replacement [S3] is Cai--Lenton--Lontzek (2016), with Judd acknowledged rather than listed as author. This correction is explicit, and intended-source identity remains uncertain. The coupled robust-control target is an editorial research family, not a conjecture asserted by those sources or certified open in 2026.

\subsection*{References}

\begin{enumerate}[label={[\textnormal{S}\arabic*]},leftmargin=*]

\item Timothy M. Lenton, Hermann Held, Elmar Kriegler, Jim W. Hall, Wolfgang Lucht, Stefan Rahmstorf and Hans Joachim Schellnhuber (2008). \emph{Tipping elements in the Earth’s climate system}. Proceedings of the National Academy of Sciences 105(6), 1786–1793. \url{https://doi.org/10.1073/pnas.0705414105}. \textbf{Locator:} Publisher or institutional abstract and bibliographic record. \textbf{Support and limits:} Physical tipping elements and expert assessment; does not identify an economic welfare distribution.

\item Derek Lemoine and Christian Traeger (2014). \emph{Watch Your Step: Optimal Policy in a Tipping Climate}. American Economic Journal: Economic Policy 6(1), 137–166. \url{https://doi.org/10.1257/pol.6.1.137}. \textbf{Locator:} Publisher or institutional abstract and bibliographic record. \textbf{Support and limits:} Specified stochastic IAM with tipping; optimal-tax claims are model conditional.

\item Yongyang Cai, Timothy M. Lenton and Thomas S. Lontzek (2016). \emph{Risk of multiple interacting tipping points should encourage rapid CO2 emission reduction}. Nature Climate Change 6, 520–525. \url{https://doi.org/10.1038/nclimate2964}. \textbf{Locator:} Publisher citation and author list, lines 206–264. \textbf{Support and limits:} Verified relevant replacement has Lenton, not Judd. Original author-year combination remains ambiguous; replacement is disclosed.

\end{enumerate}

\subsection*{Common conventions: Source mathematical conventions}

Unless an entry states otherwise, agent and firm sets are finite. State and action spaces are measurable, strategies and outcome maps are measurable, probability laws are specified, and expectations exist for the targets under discussion. Infinite-horizon discount factors lie in $(0,1)$ and discounted objectives are finite. Norms, tolerances and complexity input sizes must be specified for computational comparisons. Constants or primitives that appear locally are inputs, not empirical facts asserted by this catalogue.

\textbf{Identification.} A statistical model $m\in\mathcal M$ implies an observable law $P_m$ and target $\theta(m)$. At law $P$, its identified set is
\[
\Theta(P;\mathcal M)=\{\theta(m):m\in\mathcal M,\ P_m=P\}.
\]
Point identification holds when this set is a singleton, or equivalently when $P_{m_1}=P_{m_2}$ implies $\theta(m_1)=\theta(m_2)$ for all compatible models. Sharp bounds mean bounds attained, or approached when the boundary is not attained, by compatible models. Writing this set is a definition, not a solution: the substantive task is to establish informative restrictions and derive or estimate the set. An unrestricted-confounding model may give only trivial bounds.

\textbf{Inference.} Identification, estimability and valid uncertainty quantification are separate questions. Fix $\alpha\in(0,1)$. A confidence procedure $C_n$ over a declared law class $\mathcal P$ has asymptotic uniform coverage at level $1-\alpha$ when
\[
\liminf_{n\to\infty}\inf_{P\in\mathcal P}\Pr_P\{\theta(P)\in C_n\}\geq1-\alpha.
\]
For set-identified targets the coverage event must be defined separately, for example coverage of the full identified set. If an entry uses identification-indexed models instead of uniquely indexed laws, the same coverage convention is applied over those models. Pointwise limits do not establish uniform validity, and asymptotic validity does not guarantee finite-sample precision. Sample sizes, dependence structures and weak-identification sequences are part of each application.

\textbf{Counterfactuals.} $Y_i(a)$ is a potential outcome under a specified intervention, including its timing, implementation and versions. It is not $Y_i$ conditional on voluntarily choosing $a$. A full intervention vector permits interference; shorthand scalar treatment notation does not silently impose no interference. Assignment, selection, attrition, anticipation and measurement must be specified. A claim about an instrument's local effect is not automatically a population policy effect.

\textbf{Economic equilibrium.} A counterfactual model includes best responses, constraints, feasibility, market clearing, state transitions and expectations consistent with the stated information. When several equilibria exist, either a declared selector is an input or the target is set-valued. Comparisons across policy regimes must use the stated selection convention. Reduced-form relationships are not presumed invariant after an equilibrium-changing intervention.

\textbf{Optimization and welfare.} Optimization is over the stated feasible set. If attainment is not established, $\sup$ or $\inf$ and $\varepsilon$-optimal choices replace an asserted maximizer or minimizer. For example, with value $V=\sup_{a\in\mathcal A}W(a)<\infty$, an $\varepsilon$-optimal set is $\{a\in\mathcal A:W(a)\geq V-\varepsilon\}$ for $\varepsilon>0$. Benefits, costs, risk and health or environmental outcomes can be aggregated only using declared units and conversion weights. Interpersonal utility weights, the treatment of fiscal transfers and foreign profits, discounting, and the behavioral welfare criterion are explicit inputs. A comparative-static derivative additionally requires existence and appropriate differentiability of the selected counterfactual.

\textbf{Sources.} Local labels [S1], [S2], and so forth restart in every question. Locators identify the relevant abstract, model, section or result at the level actually checked; a bibliographic or abstract-level verification is not represented as inspection of an entire proof. Working papers are distinguished from later publications and changing versions. Source summaries are paraphrased. All URL checks and editorial formulations refer to the audit date stated above.
% END ECONOMICS STATEMENT
\end{document}
